\فصل{پیشنهادات و اهداف پیش‌رو}
\label{conclusion}
در این گزارش،  در ابتدا به تعریف قضیه‌ی جداسازی کور منابع در حالت غیرخطی پرداخته و قضایای جدایی‌پذیری موجود را بیان کردیم. در ادامه، در مورد الگوریتم‌های مبتنی بر کمینه‌کردن اطلاعات متقابل بحث کرده و تعمیمی از الگوریتم‌های موجود برای مخلوط‌های فرآیند‌های تصادفی ارائه نمودیم و عملکرد آن را بررسی کردیم. در آخر نیز به معرفی معیار 
\lr{HSIC}
به عنوان جایگزینی برای اطلاعات متقابل پرداختیم. 

قضایای جدایی‌پذیری مخلوط‌های غیرخطی، قضایایی بسیار جدید هستند و تا کنون به طور دقیق برسی نشده‌اند. این قضایا، در برخی موارد شرایط بسیار محدود کننده‌ای بر مخلوط‌ها اعمال می‌کنند و ممکن است در عمل قابل استفاده نباشند. یکی از اهداف این پروژه، بررسی دقیق شرایط این قضایا و مشخص نمودن دقیق محدویت‌هایی است که فرضیات آن‌ها دارند. حتّی با وجود این قضایا، هنوز  مسئله  اصلی ما که آیا تابعی که مشاهدات را به فرآیند‌های تصادفی مستقل نگاشت می‌کند، جداسازی کور منابع را انجام داده است یا خیر، مسئله‌ای باز است. استفاده از این قضایا برای اثبات یا رد چنین گزاره‌ای هدف اصلی این پروژه خواهد بود.

در قضایای موجود، به نظر می‌رسد وجود «نویز» در منابع ضروری باشد و ظاهراً هیچ‌کدام از آن‌ها برای منابع یقینی قابل استفاده نیستند. بررسی جامع این ادعا دشوار است و ما در پروژه‌ی کارشناسی ۲ به این مسئله خواهیم پرداخت.

قضیه‌ی 
\eqref{eq:p_tau}
در 
\cite{TCL}،
فقط برای منابعی با 
$V=1$
اثبات شده است، اما ظاهراً اثبات این قضیه برای حالت کلّی با 
$V$
دلخواه، با ایده‌هایی مشابه ایده‌ی اثبات
\eqref{eq:p_tau}
دشوار نباشد. اثبات این قضیه برای
$V$
کلّی از اهمیت بالایی برخوردار است و یکی از مسائلی است که ما در ادامه به آن خواهیم پرداخت. همچنین در اثبات مقاله‌ی 
\cite{PCL}
برای قضیه‌ی جدایی‌پذیری خود، ایرادی وجود دارد که رفع آن و ارائه‌ی اثباتی دقیق و محکم، یکی از اهداف این پروژه است. 

الگوریتم مطرح شده در فصل 
\eqref{algs}
بر روی داده‌ی مصنوعی، عملکرد بسیار مناسبی داشت. در پروژه‌ی کارشناسی دو، عملکرد این الگوریتم با الگوریتم‌های مشابه مقایسه خواهد گردید و همچنین کیفیت عملکرد آن برای حل مسئله‌ای واقعی بررسی خواهد شد. 

از آن‌جایی که برای اطلاعات متقابل دو متغیر تصادفی، تخمین‌گر مناسبی وجود ندارد، به نظر می‌رسد  با استفاده از معیار‌های دیگری مانند معیار 
\lr{HSIC}
و یا استفاده از فاصله‌های دیگر بین توزیع‌ها، قادر باشیم الگوریتم‌های بهتری برای جداسازی کور منابع، به خصوص در حالت غیرخطی، ارائه دهیم. این موضوع نیز از اهداف این پروژه است.





